\documentclass[11pt]{article}
\usepackage[margin=0.85in]{geometry}
\usepackage{booktabs}
\usepackage{array}
\usepackage{xcolor}
\usepackage{hyperref}
\usepackage{parskip}
\hypersetup{colorlinks=true, linkcolor=blue, urlcolor=blue}
\pagenumbering{gobble}

\begin{document}

\begin{center}
{\Large \textbf{free5GC Result: Frozen vs.\ Upgraded Report --- Decision Summary}}\\[0.3em]
{\small Agentic AI for Vulnerability Confirmation and Patch Validation in Open-Source Software}
\end{center}

\vspace{0.5em}

\textbf{Context.} The FYP report was frozen after several rounds of review, with the free5GC case study runtime-confirming one of four reachable handlers (the DELETE subscription handler, CVE-2026-40248) against a minimal harness. Afterward, on a separate branch, that result was extended to a fuller Docker deployment (real MongoDB, real free5GC NRF) and all four reachable handlers were tested. This surfaced a materially stronger finding: the same one-line root cause (a missing \texttt{return} after an error response) affects not just DELETE but also both GET variants and PUT, producing a full CRUD bypass rather than a single-record issue. This is now written up as an alternate, upgraded version of the report on that branch, not yet merged.

\vspace{0.5em}

\begin{center}
\renewcommand{\arraystretch}{1.4}
\begin{tabular}{@{}p{2.7cm}p{5.7cm}p{5.7cm}@{}}
\toprule
\textbf{Version} & \textbf{Strength} & \textbf{Risk / what must be accepted} \\
\midrule
\textbf{Frozen} (\texttt{main}, commit \texttt{c58483e}) &
Stable; already reviewed across multiple rounds; conservative claim. &
Only the DELETE handler is runtime-confirmed; the other three reachable handlers remain compile-verified only. \\
\addlinespace
\textbf{Upgraded} (\texttt{free5gc-full-deployment}, commit \texttt{9c95480}) &
All four reachable handlers runtime-confirmed; stronger CRUD-bypass finding (unauthorized read of one record, unauthorized read of the entire collection, unauthorized write, unauthorized delete) all from the same root cause and the same upstream fix; free5GC case becomes a complete, not partial, real-system validation. &
Newer, not yet independently reviewed; OAuth2 enforcement was explicitly disabled on the real NRF to make the handlers reachable, and this limitation must be accepted as stated, not overlooked. \\
\bottomrule
\end{tabular}
\end{center}

\vspace{0.5em}

\textbf{What changed technically (for reference).} Using the official \texttt{free5gc-compose} Docker stack (real, official MongoDB and NRF images) instead of the original minimal stand-in NRF, all three untested handlers were driven through the same malicious/benign request pattern already used for DELETE. Two were found more severe than expected: the collection-level GET returns what looks like a plain \texttt{400} rejection but the raw response body is followed by a JSON dump of every stored subscription record; the single PUT returns \texttt{404} but the submitted data is genuinely written to storage. All four close cleanly on the patched build, with legitimate use preserved. Full evidence (\texttt{manifest.json}, \texttt{verdict.json}, raw request/response logs, container logs) is committed under \texttt{free5gc\_full\_deployment/evidence/} on the branch.

\vspace{0.5em}

\textbf{Decision needed.} Should the report use the conservative, already-reviewed single-handler result (\texttt{main}), or the stronger, newer all-four-handler result (\texttt{free5gc-full-deployment})? If the latter, the upgraded PDF becomes the new report base and the frozen version is kept as the prior stable draft. If the former, the all-four-handler result is instead framed as post-report future work.

\vspace{0.5em}

\textit{Prepared as a decision aid; both report versions and the full evidence bundle are available on request.}

\end{document}
